\section{Lieb concavity theorem}

The Lieb concavity theorem rests on a scalar integral identity that expresses the
geometric interpolation $a^s b^{1-s}$ as a positive average of the elementary
fractions $ab/(a+tb)$. This is the analytic core of the integral representation
of $A^s B^{1-s}$.

\begin{lemma}[Reflection integral]\label{lem:reflection_integral}
    \lean{Real.integral_rpow_div_one_add}
    \leanok
    For $s\in(0,1)$,
    \begin{align}
        \int_0^\infty \frac{u^{s-1}}{1+u}\,du
        &= \frac{\pi}{\sin(\pi s)}.
        \notag
    \end{align}
\end{lemma}

\begin{proof}\leanok
    This is the value $B(s,1-s)$ of the Beta function, which equals
    $\Gamma(s)\Gamma(1-s)=\pi/\sin(\pi s)$ by Euler's reflection formula. The
    substitution $x=u/(1+u)$ carries the half-line onto the unit interval and
    turns the integrand into the Beta integrand $x^{s-1}(1-x)^{-s}$.
\end{proof}

\begin{lemma}[Integrability of the Lieb integrand]\label{lem:lieb_integrand_integrable}
    \lean{Real.integrableOn_lieb_integrand}
    \leanok
    For $a,b>0$ and $s\in(0,1)$, the function
    $t\mapsto t^{s-1} \dfrac{ab}{a+tb}$ is integrable on $(0,\infty)$.
\end{lemma}

\begin{proof}\leanok
    The integrand is $O(t^{s-1})$ as $t\to0^+$ and $O(t^{s-2})$ as
    $t\to\infty$, so it has integrable comparison functions on each tail.
\end{proof}

\begin{lemma}[Value of the Lieb integral]\label{lem:lieb_integral_value}
    \lean{Real.integral_lieb_integrand}
    \leanok
    \uses{lem:reflection_integral}
    For $a,b>0$ and $s\in(0,1)$,
    \begin{align}
        \int_0^\infty t^{s-1} \frac{ab}{a+tb}\,dt
        &= a^s b^{1-s} \frac{\pi}{\sin(\pi s)}.
        \notag
    \end{align}
\end{lemma}

\begin{proof}\leanok
    The scaling substitution $t=(a/b)\,u$ turns the integrand into
    $a^s b^{1-s}\,u^{s-1}/(1+u)$, whose half-line integral is the reflection
    integral $\pi/\sin(\pi s)$ of Lemma~\ref{lem:reflection_integral}.
\end{proof}

\begin{lemma}[Scalar Lieb integral identity]\label{lem:scalar_lieb_integral}
    \lean{rpow_mul_rpow_one_sub_eq_integral}
    \leanok
    \uses{lem:lieb_integral_value}
    For $a,b>0$ and $s\in(0,1)$,
    \begin{align}
        a^s b^{1-s}
        &= \frac{\sin(\pi s)}{\pi}
            \int_0^\infty t^{s-1} \frac{ab}{a+tb}\,dt.
        \notag
    \end{align}
\end{lemma}

\begin{proof}\leanok
    By Lemma~\ref{lem:lieb_integral_value}, the integral equals
    $a^s b^{1-s} \pi/\sin(\pi s)$, so
    \begin{align}
        \frac{\sin(\pi s)}{\pi}
        \int_0^\infty t^{s-1} \frac{ab}{a+tb}\,dt
        &= \frac{\sin(\pi s)}{\pi}\cdot
           a^s b^{1-s} \frac{\pi}{\sin(\pi s)}
         = a^s b^{1-s}.
        \notag
    \end{align}
\end{proof}

\begin{lemma}[Diagonal Lieb integral identity]\label{lem:diagonal_lieb_integral}
    \lean{Matrix.diagonal_lieb_integral}
    \leanok
    \uses{lem:scalar_lieb_integral}
    For positive diagonal matrices $M$ and $N$ and $s\in(0,1)$,
    \begin{align}
        M^s N^{1-s}
        &= \frac{\sin(\pi s)}{\pi}
           \int_0^\infty t^{s-1}\,M\,(M+tN)^{-1}\,N\,dt.
        \notag
    \end{align}
\end{lemma}

\begin{proof}\leanok
    \uses{lem:scalar_lieb_integral}
    Both sides are diagonal, so the identity holds entry by entry; on each
    diagonal entry it is the scalar Lieb integral identity
    (Lemma~\ref{lem:scalar_lieb_integral}).
\end{proof}

\begin{lemma}[Operator integral representation]\label{lem:op_lieb_integral_rep}
    \lean{superop_lieb_integral_rep}
    \leanok
    \uses{lem:diagonal_lieb_integral}
    For positive-definite matrices $A$, $B$ and $s\in(0,1)$,
    \begin{align}
        (A\otimes\Id)^s\,(\Id\otimes B^\top)^{1-s}
        &=
        \frac{\sin(\pi s)}{\pi}
        \int_0^\infty t^{s-1}\,(A\otimes\Id)
        ((A\otimes\Id)+t(\Id\otimes B^\top))^{-1}
        (\Id\otimes B^\top)\,dt,
        \label{eq:operator_convexity_integral_rep}
    \end{align}
    where $A\otimes\Id$ and $\Id\otimes B^\top$ are regarded as operators on
    the Kronecker model space $\C^{D\times D}$.
\end{lemma}

\begin{proof}\leanok
    The pair $(A\otimes\Id,\Id\otimes B^\top)$ is simultaneously diagonalized
    by $V=U_A\otimes U_{B^\top}$ (eigenbases of $A$ and $B^\top$).
    In that basis the pair becomes a commuting pair of positive diagonal
    matrices, and the identity follows from the diagonal case
    (Lemma~\ref{lem:diagonal_lieb_integral}).
\end{proof}

\begin{lemma}[Concavity of the parallel sum]\label{lem:parallel_sum_concave}
    \lean{Matrix.PosDef.parallel_sum_concave}
    \leanok
    For positive-definite matrices $X_1,X_2,Y_1,Y_2$ and
    $\theta\in[0,1]$, the parallel sum
    $(X,Y)\mapsto X(X+Y)^{-1}Y$ is jointly concave in the Loewner order:
    \begin{align}
        \theta\,X_1(X_1+Y_1)^{-1}Y_1
        +(1-\theta)\,X_2(X_2+Y_2)^{-1}Y_2
        &\le
        X_\theta(X_\theta+Y_\theta)^{-1}Y_\theta,
        \notag
    \end{align}
    where $X_\theta=\theta X_1+(1-\theta)X_2$ and
    $Y_\theta=\theta Y_1+(1-\theta)Y_2$.
\end{lemma}

\begin{proof}\leanok
    Write $Z=X(X+Y)^{-1}Y$. The identity
    $X(X+Y)^{-1}(X+Y)=X$ gives
    $(X-Z)-X(X+Y)^{-1}X=0$, so the block matrix
    $\begin{psmallmatrix}X-Z & X\\X & X+Y\end{psmallmatrix}$ is positive
    semidefinite (its Schur complement against $X+Y$ vanishes).
    Setting $Z_i=X_i(X_i+Y_i)^{-1}Y_i$ and
    $Z_\theta=\theta Z_1+(1-\theta)Z_2$, the convex combination of the two
    block matrices is
    $\begin{psmallmatrix}X_\theta-Z_\theta & X_\theta\\
    X_\theta & X_\theta+Y_\theta\end{psmallmatrix}$, which is positive
    semidefinite, and its Schur complement against $X_\theta+Y_\theta$ yields
    $X_\theta(X_\theta+Y_\theta)^{-1}Y_\theta-Z_\theta\ge0$, which is the
    claimed inequality (Ando, 1979).
\end{proof}

\begin{lemma}[Concavity of the resolvent integrand]
    \label{lem:superop_resolvent_integrand_concave}
    \lean{superop_resolvent_integrand_concave}
    \leanok
    \uses{lem:parallel_sum_concave}
    For $t>0$, positive-definite $A_1,A_2,B_1,B_2$, and
    $\theta\in[0,1]$, with $\hat{A}=A\otimes\Id$ and
    $\hat{B}=\Id\otimes B^\top$, setting
    $A_\theta=\theta A_1+(1-\theta)A_2$ and
    $B_\theta=\theta B_1+(1-\theta)B_2$:
    \begin{align}
        \theta\,\hat A_1(\hat A_1+t\hat B_1)^{-1}\hat B_1
        +(1-\theta) \hat A_2(\hat A_2+t\hat B_2)^{-1}\hat B_2
        &\le
        \hat A_\theta(\hat A_\theta+t\hat B_\theta)^{-1}\hat B_\theta.
        \notag
    \end{align}
\end{lemma}

\begin{proof}\leanok
    The integrand equals
    $t^{-1} \hat A(\hat A+t\hat B)^{-1}(t\hat B)$, a rescaled parallel sum of
    $\hat A$ and $t\hat B$, so concavity follows from
    Lemma~\ref{lem:parallel_sum_concave} and the linearity of
    $A\mapsto\hat A$, $B\mapsto\hat B$.
\end{proof}

\begin{lemma}[Concavity of the fractional product]\label{lem:op_lieb_concave}
    \lean{superop_lieb_concave}
    \leanok
    \uses{lem:op_lieb_integral_rep, lem:superop_resolvent_integrand_concave}
    For positive-definite matrices $A_1,A_2,B_1,B_2$, $s\in(0,1)$, and
    $\theta\in[0,1]$, with $\hat{A}=A\otimes\Id$ and
    $\hat{B}=\Id\otimes B^\top$, setting
    $A_\theta=\theta A_1+(1-\theta)A_2$ and
    $B_\theta=\theta B_1+(1-\theta)B_2$, the fractional product
    $\hat A^s\hat B^{1-s}$ is jointly concave in the Loewner order:
    \begin{align}
        \theta\,\hat A_1^{\,s}\hat B_1^{\,1-s}
        +(1-\theta) \hat A_2^{\,s}\hat B_2^{\,1-s}
        &\le
        \hat A_\theta^{\,s}\hat B_\theta^{\,1-s}.
        \notag
    \end{align}
\end{lemma}

\begin{proof}\leanok
    By the integral representation (\ref{eq:operator_convexity_integral_rep}),
    each fractional product equals the positive multiple
    $\frac{\sin(\pi s)}{\pi}$ of the integral of the resolvent integrand
    against the non-negative weight $t^{s-1}$ on $(0,\infty)$. For each $t>0$,
    the integrand is jointly concave by
    Lemma~\ref{lem:superop_resolvent_integrand_concave}, so the difference of
    the averaged point and the convex combination is, almost everywhere, the
    non-negative scalar $t^{s-1}$ times a positive-semidefinite matrix.
    Integrating a positive-semidefinite integrand yields a positive-semidefinite
    matrix, and scaling by the positive constant
    $\frac{\sin(\pi s)}{\pi}$ preserves the Loewner order, which gives the
    inequality.
\end{proof}

\begin{theorem}[Lieb concavity]\label{thm:lieb_concavity_axiom}
    \lean{lieb_concavity_axiom}
    \leanok
    For $s\in[0,1]$, any matrix~$K$, and positive-definite matrices
    $A_1,A_2,B_1,B_2$, write
    $F_s(A,B):=\Re\tr(K^\dagger A^s K B^{1-s})$. This map is jointly
    concave. For $t\in[0,1]$, set
    $A_t:=tA_1+(1-t)A_2$ and $B_t:=tB_1+(1-t)B_2$. Then
    \begin{align}
        tF_s(A_1,B_1)+(1-t)F_s(A_2,B_2)
        &\le F_s(A_t,B_t).
        \label{eq:opconv_lieb_joint_concavity}
    \end{align}
    See \cite{Lieb1973Convex,Ando1979Concavity}.
    This statement is the positive-definite, boundary-exponent ($x+y=1$) case of
    the Ando--Lieb theorem \cite[Theorem~5.15]{Wolf2012Quantum}, which holds for
    positive semidefinite $A,B$ and all exponents $x,y\ge0$ with $x+y\le1$.
    % Scope restriction recorded in docs/paper-gaps/wolf_ch5_operator_jensen_lieb.tex
\end{theorem}

\begin{proof}\leanok
    \uses{lem:op_lieb_concave}
    For $s\in(0,1)$, the concavity of the fractional product of the commuting
    left and right multiplication operators (Lemma~\ref{lem:op_lieb_concave})
    gives a Loewner-order inequality between positive-semidefinite matrices on
    the model space $\C^D\otimes\C^D$. Writing that product as
    $A^s\otimes(B^\top)^{1-s}$ and pairing it with the vectorization of $K^\top$,
    \begin{align}
        \left\langle
            \opvec K^\top,
            (A^s\otimes(B^\top)^{1-s})\opvec K^\top
        \right\rangle
        &= \Re\tr(K^\dagger A^s K B^{1-s}),
        \notag
    \end{align}
    expresses the trace functional as the quadratic form of the Kronecker
    product at that vector. Applying this positive quadratic form to the
    Loewner inequality transfers the concavity to the real part of the trace.
    The endpoints $s=0$ and $s=1$ reduce the varying side to a linear function
    of the convex-combination argument, hence hold with equality.
\end{proof}

\begin{theorem}[Lieb concavity, positive-semidefinite inputs]
    \label{thm:lieb_concavity_psd}
    \lean{lieb_concavity_psd}
    \leanok
    \uses{thm:lieb_concavity_axiom}
    For $s\in[0,1]$, any matrix~$K$, and positive-\emph{semidefinite} matrices
    $A_1,A_2,B_1,B_2$, the map
    $(A,B)\mapsto\Re\tr(K^\dagger A^s K B^{1-s})$
    is jointly concave, satisfying
    (\ref{eq:opconv_lieb_joint_concavity}) of
    Theorem~\ref{thm:lieb_concavity_axiom}.
    This is the full Ando--Lieb theorem
    \cite[Theorem~5.15]{Wolf2012Quantum} on the boundary line $x+y=1$
    (with $x=s$, $y=1-s$): the positive-definiteness restriction of
    Theorem~\ref{thm:lieb_concavity_axiom} is lifted. The general
    two-exponent region is obtained in
    Theorem~\ref{thm:lieb_concavity_subboundary}.
\end{theorem}

\begin{proof}\leanok
    \uses{thm:lieb_concavity_axiom}
    For the interior exponents $s\in(0,1)$, replace each input $A_i$ by the
    positive-definite matrix $A_i+\varepsilon\Id$, and likewise for the $B_i$.
    Theorem~\ref{thm:lieb_concavity_axiom} applies to these positive-definite
    matrices for every $\varepsilon>0$, and the average of the regularized
    inputs is the regularization of the average. Letting $\varepsilon\to0^+$,
    each regularized power $(A_i+\varepsilon\Id)^s$ converges to $A_i^s$, since
    $x\mapsto x^s$ is continuous on the non-negative reals for $0<s<1$, including
    at the origin; the trace functional is continuous, so the regularized
    inequality passes to the limit. The endpoints $s=0$ and $s=1$ are linear in
    the convex-combination argument and hold with equality.
\end{proof}

\begin{theorem}[Lieb concavity, sub-boundary exponents]
    \label{thm:lieb_concavity_subboundary}
    \lean{lieb_concavity_subboundary}
    \leanok
    \uses{thm:lieb_concavity_psd}
    Let $x,y\ge0$ with $x+y\le1$. For any matrix~$K$ and
    positive-semidefinite matrices $A_1,A_2,B_1,B_2$, write
    $F_{x,y}(A,B):=\Re\tr(K^\dagger A^x K B^y)$. This map is jointly concave.
    For $t\in[0,1]$, set $A_t:=tA_1+(1-t)A_2$ and
    $B_t:=tB_1+(1-t)B_2$. Then
    \begin{align}
        tF_{x,y}(A_1,B_1)+(1-t)F_{x,y}(A_2,B_2)
        &\le F_{x,y}(A_t,B_t).
        \notag
    \end{align}
    This is Wolf's Theorem~5.15 in the finite-dimensional matrix setting.
\end{theorem}

\begin{proof}\leanok
    \uses{thm:lieb_concavity_psd}
    Put $r=x+y$. If $r=0$, both exponents vanish and the functional is constant.
    If $0<r\le1$, set $s=x/r$. Then $x=rs$ and $y=r(1-s)$, with
    $s\in[0,1]$. The functional becomes the boundary Lieb functional applied to
    the powered inputs $A^r$ and $B^r$. Operator concavity of
    $C\mapsto C^r$ gives
    \begin{align}
        tA_1^r+(1-t)A_2^r
        &\le (tA_1+(1-t)A_2)^r,
        \notag\\
        tB_1^r+(1-t)B_2^r
        &\le (tB_1+(1-t)B_2)^r.
        \notag
    \end{align}
    The boundary positive-semidefinite theorem gives concavity at the left-hand
    powered inputs, and monotonicity of
    $(P,Q)\mapsto\Re\tr(K^\dagger P^s K Q^{1-s})$ on positive-semidefinite
    arguments transports the result to the powered convex combinations.
\end{proof}

\begin{theorem}[Map form of Lieb concavity]\label{thm:lieb_concavity}
    \lean{lieb_concavity}
    \leanok
    \uses{thm:lieb_concavity_axiom}
    For $s\in[0,1]$, any matrix~$K$, and positive-definite matrices
    $A_1,A_2,B_1,B_2$, the function
    $(A,B)\mapsto\Re\tr(K^\dagger A^s K B^{1-s})$ is jointly concave,
    satisfying (\ref{eq:opconv_lieb_joint_concavity}) of
    Theorem~\ref{thm:lieb_concavity_axiom}.
\end{theorem}

\begin{proof}\leanok
    Follows directly from Theorem~\ref{thm:lieb_concavity_axiom}.
\end{proof}

\begin{corollary}[Lieb concavity, $K=\Id$ case]\label{cor:lieb_concavity_id}
    \lean{lieb_concavity_id}
    \leanok
    \uses{thm:lieb_concavity}
    For $s\in[0,1]$, $t\in[0,1]$, and positive-definite matrices
    $A_1,A_2,B_1,B_2$, the map
    $(A,B)\mapsto\Re\tr(A^s B^{1-s})$ is jointly concave:
    \begin{align}
        t\,\Re\tr(A_1^s B_1^{1-s})
        +(1-t)\Re\tr(A_2^s B_2^{1-s})
        &\le
        \Re\tr\!\bigl(        (tA_1+(1-t)A_2)^s
            (tB_1+(1-t)B_2)^{1-s}
        \bigr).
        \notag
    \end{align}
    Obtained from Theorem~\ref{thm:lieb_concavity} by taking $K=\Id$.
\end{corollary}

\begin{proof}\leanok
    Specialize Theorem~\ref{thm:lieb_concavity} at $K=\Id$ and simplify
    $\Id^\dagger=\Id$ together with $\Id\cdot X=X\cdot\Id=X$.
\end{proof}

\begin{corollary}[Concavity of Lieb's map in the first argument]
    \label{cor:lieb_concavity_in_fst}
    \lean{lieb_concavity_in_fst}
    \leanok
    \uses{thm:lieb_concavity}
    For $s\in[0,1]$, $t\in[0,1]$, matrix $K$, and positive-definite matrices
    $A_1,A_2,B$, the map
    $A\mapsto\Re\tr(K^\dagger A^s K B^{1-s})$ is concave:
    \begin{align}
        t\,\Re\tr(K^\dagger A_1^s K B^{1-s})
        +(1-t)\Re\tr(K^\dagger A_2^s K B^{1-s})
        &\le
        \Re\tr\!\bigl(        K^\dagger(tA_1+(1-t)A_2)^s K B^{1-s}
        \bigr).
        \notag
    \end{align}
\end{corollary}

\begin{proof}\leanok
    Specialize Theorem~\ref{thm:lieb_concavity} at $B_1=B_2=B$, so that
    the convex combination collapses: $tB+(1-t)B=B$.
\end{proof}

\begin{corollary}[Concavity of Lieb's map in the second argument]
    \label{cor:lieb_concavity_in_snd}
    \lean{lieb_concavity_in_snd}
    \leanok
    \uses{thm:lieb_concavity}
    For $s\in[0,1]$, $t\in[0,1]$, matrix $K$, and positive-definite matrices
    $A,B_1,B_2$, the map
    $B\mapsto\Re\tr(K^\dagger A^s K B^{1-s})$ is concave:
    \begin{align}
        t\,\Re\tr(K^\dagger A^s K B_1^{1-s})
        +(1-t)\Re\tr(K^\dagger A^s K B_2^{1-s})
        &\le
        \Re\tr\!\bigl(        K^\dagger A^s K(tB_1+(1-t)B_2)^{1-s}
        \bigr).
        \notag
    \end{align}
\end{corollary}

\begin{proof}\leanok
    Specialize Theorem~\ref{thm:lieb_concavity} at $A_1=A_2=A$, so that
    the convex combination collapses: $tA+(1-t)A=A$.
\end{proof}

\section{Resolvent monotonicity toward the Lieb concavity theorem}

The integral-representation route toward the Lieb concavity theorem rests on the
antitonicity of the resolvent of the commuting left- and right-multiplication
superoperators. Its matrix-order foundation is the inverse-antitonicity of the
Loewner order.

\begin{lemma}[Loewner inverse-antitonicity]\label{lem:loewner_inv_antitone}
    \lean{Matrix.PosDef.inv_le_inv_of_le}
    \leanok
    For positive-definite matrices $A$ and $B$ with $A\le B$ in the Loewner
    order, the inverses satisfy $B^{-1}\le A^{-1}$.
\end{lemma}

\begin{proof}\leanok
    Compare the two Schur complements of the block matrix
    $\begin{psmallmatrix}A^{-1}&\Id\\\Id&B\end{psmallmatrix}$. Complementing its
    $(1,1)$ block gives $B-(A^{-1})^{-1}=B-A\ge0$, so the block matrix is
    positive semidefinite; complementing the $(2,2)$ block of the same matrix
    gives $A^{-1}-B^{-1}$, which is therefore positive semidefinite as well.
\end{proof}

\begin{lemma}[Joint Loewner-antitonicity of the commuting-multiplication resolvent]
    \label{lem:superop_resolvent_antitone}
    \lean{superop_resolvent_antitone}
    \leanok
    \uses{lem:loewner_inv_antitone}
    Fix $t>0$ and positive-definite matrices with $A_1\le A_2$ and
    $B_1\le B_2$. The resolvent of the Kronecker model
    $A\otimes\Id+t\,(\Id\otimes B^\top)$ of the commuting left- and
    right-multiplication superoperators is antitone:
    \begin{align}
        (A_2\otimes\Id+t\,(\Id\otimes B_2^\top))^{-1}
        &\le
        (A_1\otimes\Id+t\,(\Id\otimes B_1^\top))^{-1}.
        \notag
    \end{align}
\end{lemma}

\begin{proof}\leanok
    \uses{lem:loewner_inv_antitone}
    Each operator $A\otimes\Id+t\,(\Id\otimes B^\top)$ is positive definite,
    being the sum of the positive-definite $A\otimes\Id$ and the
    positive-semidefinite $t\,(\Id\otimes B^\top)$. From $A_1\le A_2$ and
    $B_1\le B_2$ one gets $A_1\otimes\Id\le A_2\otimes\Id$ and
    $\Id\otimes B_1^\top\le\Id\otimes B_2^\top$; since $t>0$, adding these
    yields
    \begin{align}
        A_1\otimes\Id+t\,(\Id\otimes B_1^\top)
        &\le
        A_2\otimes\Id+t\,(\Id\otimes B_2^\top).
        \notag
    \end{align}
    Inverse-antitonicity (Lemma~\ref{lem:loewner_inv_antitone}) reverses this
    inequality on passing to the resolvents.
\end{proof}
